% Included within the parent's evaluation/provenance appendix; no top-level section.
\subsection{Evaluation provenance and missing outcomes}
\label{app:evaluation-provenance}
The historical R81--R83 public-source and component-ablation tables use archived source-referenced pixel measurements at $512\times512$. Whole-image NCC is distinct from the historical local-pattern NCC. With a full-image crop, only zero displacement is valid in the shift search, so fixed and shifted NCC coincide. Official PIE masks define edit regions; their complements define preservation regions. Empty complements are excluded from preservation MAE, with source and record counts reported in each table. Edit-region MAE measures pixel departure, while requested-edit achievement requires a separate outcome.

\paragraph{Record accounting and evaluation status.}
The pilot and historical-source sweep (archive R81) completed 366 of 366 outputs: 240 from five methods on 24 public sources and two seed slots, plus 126 from nine controls on seven historical tasks and two seeds. The support ablation (R82) completed 80 of 80 new outputs; sixty frozen AUTO\_DETAIL, U25\_KV, and FlowEdit N26 records yield its 140-record comparison. The memory ablation (R83) completed 60 of 60 new outputs; eighty frozen U25\_KV, FlowEdit N26, SOFT75, and REQUEST\_SOFT25 records yield its separate 140-record comparison. Each ablation table uses the same ten development sources and twenty source--seed slots per method. The support-only reference is displayed as REQUEST\_SOFT25 (memory-fixed) throughout the results.

CLIP, DINO, LPIPS, and formal human scores are absent from the historical R81--R83 comparisons, whose human-rating fields remain unfilled. Final R84/R85 automatic evaluation is complete and reported separately below. Specifically, the support ablation records eighty \texttt{pipeline\_error} states caused by a missing local CLIP snapshot; LPIPS was not requested. The memory ablation records sixty \texttt{not\_run} states, while its frozen references retain forty pilot rows marked \texttt{not\_run\_no\_native\_png\_in\_compact\_return} and forty support-ablation pipeline-error rows. These statuses concern evaluation execution; new generation and recorded pixel measurements completed. Missing outcomes remain unknown and are not converted into zero scores or failed-edit labels.

\paragraph{Resampling and timing.}
The pilot uses 20,000 percentile source-bootstrap resamples with random seed 81030; its two-sided sign-flip calculation uses 20,000 draws, plus-one correction, and Holm adjustment across four comparators within each metric. The support ablation uses 10,000 source-bootstrap resamples with seed 42. The memory ablation uses 20,000 with seed 83, resetting the same generator for each metric and contrast. Both component ablations report unadjusted exploratory intervals after averaging paired seeds within each source. Bootstrap seeds are distinct from the generation seeds 0 and 17.

Recorded times include run-level variability. SOFT25's first recorded slot takes 84.77 seconds; its mean is 53.25 seconds, versus 45.42 seconds for REQUEST\_SOFT25 (memory-fixed), despite the latter requiring six additional transformer calls. REQUEST\_MEMORY ranges from 45.26 to 77.45 seconds and averages 55.63 seconds; REQUEST\_BOTH and MEAN\_MEMORY average 45.58 and 45.80 seconds. These observations do not identify an intrinsic speed advantage; fixed NFE and recorded timing are reported separately.

\paragraph{Artifacts and interpretation.}
Pixel values in the tables are archived measurements from lossless outputs and source evaluation PNGs. Compact returns contain JPEG previews and scalar step summaries, omitting the original output PNGs and full per-job records; the frozen pixel measurements were not recomputed from JPEGs. The support and memory reviews each cover all twenty paired sheets, with selected original previews shown in the figures. Those observations document specific edits and artifacts rather than an aggregate semantic success rate. Lower source correlation can accompany a requested style change or unwanted drift; interpretation uses the requested transformation and the visible outcome in addition to source-referenced metrics.

\paragraph{RF-Edit implementation provenance.}
The adapter's time signs, half-step slots, injection intervals, and text-plus-image value replacement agree with the pinned official implementation, commit \texttt{7ac380143379b656eca5398157f41f97ef1ef2f9}. In historical R81, its 25 solver steps and 20 injection intervals follow that implementation's defaults, with 100 transformer evaluations. Final R84/R85 RF-Edit I08 instead uses eight injection intervals and one deterministic seed, retaining 100 transformer evaluations. In R81, the two seed slots reproduce the same preview for each task and are repeated execution records within a source. The recorded generic field \texttt{memory\_enabled=false} describes the internal TraceMemory hook; RF-Edit separately records 720 feature captures and 720 applications. The revised reporting interface distinguishes the two mechanisms. The adapter integrates latent states in FP32, whereas the official implementation updates them in BF16.

\paragraph{Experimental scope.}
The 24-source pilot was source-disjoint from the historical six-source development set at selection. Its images subsequently informed development; the ten-source support and memory subset is therefore a seen development set with one source per official category. Mean-gate matching constrains the average query multiplier, not attention-output energy. The small spatial-versus-mean estimates neither establish equivalence nor rule out a task-specific benefit. Area-aware normalization and restricted source-pool proposals remain untested extensions; the completed ablations change support floors, correction allocation, and query-memory gates as specified in the results.

\paragraph{Regional field names and archived analysis.}
The compact table names edit-region error \texttt{object\_mae} and preservation-region error \texttt{bg\_mae}. The sixty new rows have sixty edit-region values and forty-eight preservation-region values; the other twelve preservation entries correspond to the two empty complements. The accompanying \texttt{tables/r83\_evidence} directory includes the input metric tables, summary, source-paired estimates, means, and analysis JSON. These archived scalar values support aggregation without the original PNG files. Recomputing pixel or embedding metrics from images requires the native lossless outputs.

\subsection{Completed final R84/R85 evaluation provenance}
\label{app:final-provenance}
The released numerical evidence resides in the repository's \path{results/R84/} and \path{results/R85/} directories. The \path{analysis/} subdirectories contain source-level means, paired comparisons, and numeric summaries; \path{evidence/} retains the original automated source-bootstrap and sign-flip exports. The final metrics were evaluated from native lossless outputs. Compact image previews support visual inspection.

R85 records 520/520 successful automatic evaluations across 40 sources: six arms have 80 outputs each and RF-Edit has 40. Seed sets are complete and arm--source--seed entries are unique. R84 records 130/130 successful automatic evaluations across ten sources. Seven R85 sources and two R84 sources have empty official preservation complements; the corresponding regional metrics are omitted, leaving 33 and eight eligible sources. Their absence is structural and does not indicate failed evaluation.

The final estimand is the equal-source mean after averaging planned seeds within each arm and source. R85 uses 10,000 percentile source-bootstrap resamples and two-sided source sign-flip Monte Carlo tests with 10,000 replicates and plus-one correction. The main adjusted $p$ values use six CORR contrasts separately per metric; BOTH--CORR reverses one contrast's sign. Pointwise intervals have no multiplicity adjustment. The all-pair 21-comparison and global 66-numeric-test corrections are sensitivity analyses, with multiplicity sensitivity reported below. R84 uses 10,000 paired source-bootstrap resamples and exact source sign flips for its small population, with six planned comparisons per metric. RF-Edit's single result remains one observation per source, compared with the other arm's within-source seed mean.

CLIP uses ViT-L/14, DINO uses DINOv2-base full-image features, and LPIPS uses AlexNet. Preservation LPIPS weights its spatial map by the area-resampled official edit-mask complement; feature receptive fields may cross the boundary. The archived \texttt{pixel\_ssim} is a background-region metric whose local windows can also cross that boundary. Full-image DINO and a crop that spans the whole source do not supply an independent local-pattern measurement. NFE records complete standalone calls: CORR, BOTH, and ENERGY 97; U25 91; FlowEdit N26 52; ETA03 56; RF-Edit I08 100. R84 FlowEdit N28 uses 56 calls and remains a development control, separate from the final N26 benchmark.

\paragraph{Multiplicity sensitivity of the final comparison.}
Multiplicity sensitivity narrows the preservation result: its adjusted $p$ is 0.070193 over all 21 arm pairs within the metric and 0.285571 over the additional 66-test numeric family. The full-image LPIPS improvement remains significant under the 66-test adjustment ($p=0.023798$). The principal preservation statement uses the declared six-comparison within-metric family.
